\documentclass[11pt]{amsart}

\usepackage[T1]{fontenc}
\usepackage[utf8]{inputenc}
\usepackage{lmodern}
\usepackage{amsmath,amssymb,amsthm,mathtools}
\usepackage{enumitem}
\usepackage{microtype}
\usepackage[colorlinks=true,linkcolor=blue,citecolor=blue,urlcolor=blue]{hyperref}
\usepackage[nameinlink,capitalise]{cleveref}

\setlist[itemize]{leftmargin=2em}
\setlist[enumerate]{leftmargin=2em}

\newtheorem{theorem}{Theorem}[section]
\newtheorem{proposition}[theorem]{Proposition}
\newtheorem{corollary}[theorem]{Corollary}
\newtheorem{lemma}[theorem]{Lemma}
\newtheorem{question}[theorem]{Question}
\theoremstyle{definition}
\newtheorem{definition}[theorem]{Definition}
\newtheorem{remark}[theorem]{Remark}

\DeclareMathOperator{\disc}{disc}
\DeclareMathOperator{\discop}{disc_{op}}
\DeclareMathOperator{\diag}{diag}
\DeclareMathOperator{\spanop}{span}
\newcommand{\C}{\mathbb C}
\newcommand{\R}{\mathbb R}
\newcommand{\T}{\mathbb T}
\newcommand{\eps}{\varepsilon}
\newcommand{\op}{\mathrm{op}}
\newcommand{\norm}[1]{\left\lVert #1\right\rVert}
\newcommand{\abs}[1]{\left\lvert #1\right\rvert}
\newcommand{\ip}[2]{\left\langle #1,#2\right\rangle}

\title[Matrix Spencer for Toeplitz and Hankel matrices]{Matrix Spencer for finite Toeplitz and Hankel matrices}
\author{Yitzchak Shmalo}
\address{Einstein Institute of Mathematics, The Hebrew University of Jerusalem, Jerusalem, Israel}
\date{June 2026}

\begin{document}

\subjclass[2020]{Primary 15A60, 11K38; Secondary 15B05, 42A05, 47B35, 68Q25}
\keywords{Matrix discrepancy, Matrix Spencer conjecture, Toeplitz matrices, Hankel matrices, trigonometric polynomials, finite Toeplitz extension, vector balancing}

\begin{abstract}
The Matrix Spencer problem asks whether every family of self-adjoint
$n\times n$ contractions admits signs whose signed sum has operator norm
$O(\sqrt n)$.  We prove this bound for finite scalar Toeplitz and Hankel
matrices.  More precisely, for arbitrary complex $d\times d$ Toeplitz
contractions $A_1,\ldots,A_n$,
\[
   \discop(A_1,\ldots,A_n)
   \le C\sqrt{n\log\!\left(e+\frac dn\right)} .
\]
The same estimate holds for arbitrary complex Hankel contractions.  In
particular, self-adjoint scalar Toeplitz and Hankel matrices satisfy the
conjectured $O(\sqrt n)$ Matrix Spencer bound throughout the regime $d\le n$.
The proof is short and uses only three ingredients: the finite scalar Toeplitz
extension theorem of Nikolskaya--Farforovskaya, de la Vall\'ee Poussin
localization, and the rectangular vector form of Spencer's theorem on an
$O(d)$-point norming grid.  We also record two direct finite-section variants:
a bandwidth-$w$ estimate with entropy $w$, and a tame full-bandwidth estimate
under the scale-invariant condition $\|p_i'\|_\infty=O(d)$.
\end{abstract}

\maketitle

\tableofcontents

\section{Introduction}

Spencer's theorem that six standard deviations suffice is a central result of
combinatorial discrepancy theory \cite{Spencer1985}.  In a rectangular form, if
$x_1,\ldots,x_n$ lie in the unit ball of $\ell_\infty^m$, then there are signs
$\eps_i\in\{-1,1\}$ such that
\[
   \left\|\sum_{i=1}^n \eps_i x_i\right\|_\infty
   \le C\sqrt{n\log\!\left(e+\frac mn\right)} .
\]
This form follows from the classical entropy method; see, for example,
Gluskin \cite{Gluskin1989}, Chazelle \cite{Chazelle2000}, Matousek
\cite{Matousek1999}, and Giannopoulos \cite{Giannopoulos1997}.

The Matrix Spencer problem asks for an analogous estimate with vectors replaced
by self-adjoint matrices and the $\ell_\infty$ norm replaced by operator norm.
We use the following notation.

\begin{definition}[Matrix discrepancy]
For complex matrices $A_1,\ldots,A_n$ of the same size, define
\[
   \discop(A_1,\ldots,A_n)
   := \min_{\eps\in\{-1,1\}^n}
      \left\|\sum_{i=1}^n \eps_i A_i\right\|_{op} .
\]
When the matrices are self-adjoint we also use the shorter notation
\[
   \disc(A_1,\ldots,A_n)=\discop(A_1,\ldots,A_n).
\]
\end{definition}

\begin{question}[Matrix Spencer]
Is there a universal constant $C$ such that, for every $n$ and every family of
self-adjoint $n\times n$ contractions $A_1,\ldots,A_n$,
\[
   \disc(A_1,\ldots,A_n)\le C\sqrt n ?
\]
Equivalently for the applications below, one asks for the same estimate for
$d\times d$ self-adjoint contractions whenever $d\le n$, with $C$ independent of
$d$ and $n$.
\end{question}

The conjecture was popularized by Meka \cite{Meka2014}.  General random-sign
arguments based on noncommutative Khintchine inequalities give logarithmic
dimensional losses \cite{LustPiquardPisier1991}.  Several important cases are
known: rank-one and Kadison--Singer type settings
\cite{MSS2015,KyngLuhSong2020}, low-rank and Frobenius-norm regimes
\cite{BansalJiangMeka2024,DadushJiangReis2022}, and consequences of intrinsic
freeness and matrix concentration \cite{BBvH2023,BCSvH2026}.  These methods do
not directly cover general finite Toeplitz contractions, which can have rank
$d$ and Frobenius norm of order $\sqrt d$.

A scalar trigonometric polynomial
\[
   p(\theta)=\sum_{|k|<d} c_k e^{ik\theta}
\]
defines the finite Toeplitz section
\[
   T_d(p)=(c_{x-y})_{0\le x,y<d} .
\]
The matrix is self-adjoint exactly when $p$ is real-valued.  A Hankel matrix is
constant on anti-diagonals.  These classes are among the most basic structured
full-rank matrix families; standard background on Toeplitz matrices and
Toeplitz operators may be found in \cite{GrenanderSzego1958,Gray2006,BoettcherGrudsky2005,BoettcherSilbermann2006}.

\subsection{Main result}

\begin{theorem}[Toeplitz and Hankel Matrix Spencer]\label{thm:main}
There is a universal constant $C$ such that the following holds.  Let
$A_1,\ldots,A_n\in M_d(\C)$ be arbitrary complex Toeplitz matrices satisfying
$\norm{A_i}_{op}\le1$.  Then
\[
   \discop(A_1,\ldots,A_n)
   \le C\sqrt{n\log\!\left(e+\frac dn\right)} .
\]
The same estimate holds for arbitrary complex Hankel matrices.  Consequently,
self-adjoint scalar Toeplitz and Hankel matrices satisfy Matrix Spencer in the
regime $d\le n$.
\end{theorem}

The proof is given in \cref{sec:main-proof}.  The finite Toeplitz extension
theorem supplies, for each finite Toeplitz contraction $A_i$, a bounded
$L^\infty$ symbol with the same Fourier coefficients in the range seen by
$T_d$.  Convolution with a de la Vall\'ee Poussin kernel turns this symbol into a
trigonometric polynomial of degree $O(d)$ without changing the finite Toeplitz
section.  An $O(d)$-point grid norms every such polynomial up to an absolute
constant, so rectangular vector Spencer can be applied to the sampled values.
The signed polynomial is then transferred back to operator norm by the elementary
Toeplitz compression estimate $\norm{T_d(f)}_{op}\le\norm{f}_\infty$.

\subsection{Finite-section refinements}

The main theorem uses a global bounded extension theorem, and therefore its
entropy parameter is $d$.  Direct finite-section arguments give sharper or more
intrinsic statements in two common situations.

\begin{theorem}[Banded Toeplitz signing]\label{thm:banded-intro}
There is a universal constant $C$ such that the following holds.  Let
$w\ge1$ and $d\ge20w$.  Let $A_1,\ldots,A_n\in M_d(\C)$ be Toeplitz matrices of
bandwidth at most $w$, and assume $\norm{A_i}_{op}\le1$ for all $i$.  Then
\[
   \discop(A_1,\ldots,A_n)
   \le C\sqrt{n\log\!\left(e+\frac wn\right)} .
\]
In particular, if $w\le n$, then the discrepancy is $O(\sqrt n)$.
\end{theorem}

\begin{theorem}[Tame full-bandwidth Toeplitz signing]\label{thm:tame-intro}
There is a universal constant $C$ such that the following holds.  Let
$A_i=T_d(p_i)$, where each $p_i$ is a complex trigonometric polynomial of degree
less than $d$, and assume $\norm{A_i}_{op}\le1$.  Suppose
\[
   \max_{1\le i\le n}\norm{p_i'}_\infty\le L .
\]
Then
\[
   \discop(A_1,\ldots,A_n)
   \le C\left(1+\frac Ld\right)
      \sqrt{n\log\!\left(e+\frac dn\right)} .
\]
In particular, if $d\le n$ and $L=O(d)$, then the discrepancy is $O(\sqrt n)$.
\end{theorem}

The proofs of these refinements are included in \cref{app:finite-section}.  They
are not needed for \cref{thm:main}.

\subsection{Scope and relation to recent algebraic work}

This paper is focused on scalar Toeplitz and Hankel matrices.  A broader
finite-block lifting language is useful for organizing related examples, and a
fixed-block version is recorded in \cref{app:lifts}.  We keep that material in
an appendix because it is not needed for the scalar theorem and because recent
work of Akbas--Sra \cite{AkbasSra2026} proves an algebraic Matrix Spencer
theorem for finite-dimensional $C^*$-algebras, while Bandeira--B{\"o}lcskei
\cite{BandeiraBolcskei2026} proves the Matrix Spencer bound for the signed list
of left-regular representation matrices of a finite group.  The new structured
case treated here is the scalar finite Toeplitz/Hankel family, obtained through
bounded Toeplitz extension rather than through a common finite-dimensional
matrix algebra.

\section{Preliminaries}\label{sec:prelim}

We identify $\T$ with $\R/(2\pi\mathbb Z)$ and use normalized Haar measure
$d\theta/(2\pi)$.  For $f\in L^1(\T)$, write
\[
   \widehat f(k)=\int_{\T} f(\theta)e^{-ik\theta}\,\frac{d\theta}{2\pi} .
\]
If $f\in L^\infty(\T)$, the $d\times d$ Toeplitz section induced by $f$ is
\[
   T_d(f)=(\widehat f(x-y))_{0\le x,y<d} .
\]
When $p$ is a trigonometric polynomial of degree less than $d$, this agrees with
the finite Toeplitz matrix determined by the Fourier coefficients of $p$ in the
range $|k|<d$.

\subsection{Vector Spencer}

\begin{theorem}[Rectangular vector Spencer]\label{thm:vector-spencer}
There is a universal constant $C_S$ such that, for every
$x_1,\ldots,x_n\in\C^m$ with $\norm{x_i}_\infty\le1$, there are signs
$\eps\in\{-1,1\}^n$ satisfying
\[
   \left\|\sum_{i=1}^n \eps_i x_i\right\|_\infty
   \le C_S\sqrt{n\log\!\left(e+\frac mn\right)} .
\]
\end{theorem}

For complex vectors this follows by applying the real theorem to the real and
imaginary parts and absorbing the resulting absolute factor.  The real theorem
is the standard rectangular form of Spencer's theorem
\cite{Spencer1985,Gluskin1989,Chazelle2000,Matousek1999,Giannopoulos1997}.

\subsection{Toeplitz compression}

\begin{lemma}[Toeplitz upper bound]\label{lem:toeplitz-upper}
For every $f\in L^\infty(\T)$ and every $d\ge1$,
\[
   \norm{T_d(f)}_{op}\le \norm{f}_\infty .
\]
If $f$ is real-valued, then $T_d(f)$ is self-adjoint.
\end{lemma}

\begin{proof}
For $z,y\in\C^d$, set
$q_z(\theta)=\sum_{x=0}^{d-1}z_x e^{ix\theta}$ and define $q_y$ similarly.  Then
\[
   \ip{T_d(f)z}{y}
   =\int_\T f(\theta)q_z(\theta)\overline{q_y(\theta)}\,
       \frac{d\theta}{2\pi} .
\]
Cauchy's inequality and Parseval give
$\abs{\ip{T_d(f)z}{y}}\le\norm{f}_\infty\norm z_2\norm y_2$.  Taking the
supremum over unit vectors proves the norm bound.  If $f$ is real-valued, then
$\widehat f(-k)=\overline{\widehat f(k)}$, so $T_d(f)$ is self-adjoint.
\end{proof}

\subsection{Norming grids}

We shall use Bernstein's inequality for trigonometric polynomials
\cite[Chapter X]{Zygmund1959}.

\begin{lemma}[Bernstein inequality]\label{lem:bernstein}
If $P$ is a trigonometric polynomial of degree at most $N$, then
\[
   \norm{P'}_\infty\le N\norm{P}_\infty .
\]
\end{lemma}

\begin{lemma}[Scalar norming grid]\label{lem:grid}
Let $P$ be a complex trigonometric polynomial of degree at most $N$.  Let
$Q\subset\T$ be a uniform grid of size at least $8N$ when $N\ge1$, and let
$Q$ consist of one point when $N=0$.  Then
\[
   \norm{P}_\infty\le2\max_{\theta\in Q}\abs{P(\theta)} .
\]
\end{lemma}

\begin{proof}
Assume $N\ge1$, since the constant case is immediate.  Choose $\theta_*$ with
$\abs{P(\theta_*)}=\norm{P}_\infty$, and choose $\theta\in Q$ with
$\abs{\theta-\theta_*}\le\pi/|Q|$.  By \cref{lem:bernstein},
\[
   \abs{P(\theta_*)-P(\theta)}
   \le N\norm{P}_\infty\frac{\pi}{|Q|}
   \le\frac\pi8\norm{P}_\infty .
\]
Thus
$\max_Q\abs P\ge(1-\pi/8)\norm{P}_\infty\ge\frac12\norm{P}_\infty$.
\end{proof}

\subsection{De la Vallee Poussin kernels}

Let $F_M$ denote the Fejer kernel with Fourier coefficients
\[
   \widehat F_M(k)=\max\!\left(0,1-\frac{|k|}{M}\right),
   \qquad M\ge1 .
\]
Thus $F_M\ge0$ and $\norm{F_M}_{L^1}=1$.  Define
\[
   V_M=2F_{2M}-F_M .
\]
Then
\[
   \widehat V_M(k)=1 \quad (|k|\le M),
   \qquad
   \widehat V_M(k)=0 \quad (|k|\ge2M),
\]
and
\[
   \norm{V_M}_{L^1}\le3 .
\]
The strict inequalities at the endpoints are immaterial for us; we only use the
